% ---
% id: statistics/binomial-box/req:rule-in-force
% kind: req
% title: "The rule the request selects"
% summary: >-
%   The engine builds the box with the endpoint rule and tail evaluation the
%   request selects: the exact rule evaluated by its closed forms, which is the
%   default, the enclosure evaluation of the exact rule, or the Bernstein rule.
%   It runs no other rule.
% prominence: standard
% status: implemented
% uses: [def:endpoint-rule, def:interval]
% mentions: [prop:rule-exact, prop:rule-bernstein, prop:bracketed-endpoints, def:construction, req:completeness-report]
% assumes: []
% refs: []
% provenance: newly-derived
% implements: certiqs_key.statistics.rules:resolve_rule
% lean:
% ---

Each rule of this submodule is an \tref{admissible endpoint rule}, so every
result stated for all such rules holds for the box the engine builds,
whichever the request chooses; a result that needs more of the rule, such as
monotonicity in the allocation, holds on the rules whose propositions prove
it.

\begin{requirement}\label{req:rule-in-force}
The engine builds the box of Definition~\ref{def:construction} with the rule
that the request's \iface{endpoint_rule} and \iface{tail_evaluation} select,
and with no other: \iface{endpoint_rule: exact} with
\iface{tail_evaluation: sweep}, the \tref{exact binomial quantile rule} of
Definition~\ref{def:interval} with its tail comparisons decided on the
closed forms, exactly or against upper bounds on the tails as
Proposition~\ref{prop:rule-exact}\,(vii) admits, which is the default; \iface{endpoint_rule: exact} with
\iface{tail_evaluation: enclosure}, the enclosure evaluation
$R_{\mathrm{encl}}$ of the exact rule
(Proposition~\ref{prop:bracketed-endpoints}), an admissible relaxation of it;
and \iface{endpoint_rule: bernstein}, the Bernstein rule of
Proposition~\ref{prop:rule-bernstein}, on which \iface{tail_evaluation}
selects only how the completeness figure is computed
(Requirement~\ref{req:completeness-report}).
\end{requirement}
